\twocolumn[{
\section{Requirement-to-Evidence Map}\label{app:map}
This map concerns the written report. The assignment's separate slide,
oral-presentation, balanced-participation, and Q\&A requirements require
presentation deliverables and cannot be fulfilled by a report alone.
The main text occupies pages 1--\pageref{LastBodyPage};
references and appendices are separate.

\begin{center}
\begin{minipage}{\textwidth}
\centering\small
\captionsetup{hypcap=false}
\captionof{table}{Coverage of the supplied CS420 S06 seminar brief.}
\begin{tabularx}{\linewidth}{@{}p{1.1cm}YY@{}}
\toprule
Task & Required content & Location / evidence\\
\midrule
1 & Generative versus discriminative; GAN, autoregressive, diffusion &
Sections~\ref{sec:families}, \ref{sec:generator-comparison};
Table~\ref{tab:families}\\
2 & Forward process, beta schedule, Gaussian noise, noise predictor,
reverse sampling; both-direction diagram &
Section~\ref{sec:ddpm}; Eqs.~\eqref{eq:forward}--\eqref{eq:reverse};
Figure~\ref{fig:ddpm}\\
3 & Manual 1D forward step, symbols, stochastic expectation &
Section~\ref{sec:scalar}; $x_1=0.600833$; conditional mean and variance\\
4 & Small image at five noise levels; optional pretrained demonstration &
Section~\ref{sec:noise}; Figure~\ref{fig:noise}; Table~\ref{tab:noise};
executable code and arrays. Pretrained generation was not performed\\
5 & VAE encoder/decoder, latent operations, cross-attention, pipeline &
Sections~\ref{sec:vae}--\ref{sec:attention};
Eqs.~\eqref{eq:vae}--\eqref{eq:attention};
Figures~\ref{fig:ldm}--\ref{fig:attention}\\
6 & Speed, quality, prompt control, data, copyright, bias;
explicit pixel/latent comparison &
Sections~\ref{sec:cfg}, \ref{sec:benchmark}--\ref{sec:ethics};
Tables~\ref{tab:benchmark}--\ref{tab:comparison}\\
\bottomrule
\end{tabularx}
\end{minipage}
\end{center}
\vspace{8pt}
}]


\textbf{Required representations.}
Equation~\eqref{eq:forward} is the Gaussian forward transition;
Eq.~\eqref{eq:loss} is the simplified noise-prediction objective.
The worked example gives a complete substitution, result, and interpretation.
Appendix~\ref{app:notation} records ranges, units, and modeling assumptions.

\textbf{Paper reading and synthesis.}
DDPM's forward/reverse method and contribution are reviewed in
Section~\ref{sec:ddpm}. LDM's representation, method, conditioning,
contributions, and differences are developed in
Sections~\ref{sec:ldm}--\ref{sec:synthesis}.
Section~\ref{sec:benchmark} identifies the exact source table for performance
claims. The report includes Abstract, Introduction, Background/Paper Review,
Method/Worked Example, Results and Discussion, Conclusion, and References.
English is used according to the commissioning instruction, which overrides
the brief's Vietnamese-language request.

\pagecut
\section{Reproducibility, Evaluation Protocol, and Attribution}\label{app:repro}
\subsection{Reproduce the original illustration}
The Overleaf project already contains the generated figure; Python is not
needed to compile the report. To regenerate it locally:
\begin{lstlisting}[language=bash]
python -m pip install -r code/requirements.txt
python code/make_figures.py
latexmk -pdf -interaction=nonstopmode -halt-on-error main.tex
\end{lstlisting}
The figure uses a synthetic geometric image, not MNIST or another dataset.
The script records the seed, schedule coefficients, library versions,
scalar example, and benchmark-ratio arithmetic in
\texttt{figures/analysis.json}. It saves unclipped frames, timesteps, the
schedule, and cumulative coefficients in \texttt{noise\_arrays.npz},
and summary coefficients in \texttt{noise\_levels.csv}.
The delivered run used NumPy 2.3.5 and ReportLab 4.4.9.
The sequence is one sampled Markov path; identical seeds and library
versions reproduce it. No image-generation checkpoint is involved.

\subsection{A future prompt and runtime study (not executed)}
Select and record exact checkpoint identifiers, revisions, dataset licenses,
hardware, software versions, numerical precision, resolution, and batch
size. Keep each fixed within a comparison and record warm-up policy.
For latent models, measure text encoding, denoising, decoding, and
end-to-end latency separately, and record peak allocated memory over the
entire pipeline.

Use paired prompts such as \emph{a red cube to the left of a blue sphere}
and \emph{a green cube to the left of a blue sphere}; separately swap the
spatial relation. Predefine at least 20 seeds per condition and inspect
all samples. Vary sampling steps and guidance in separate sweeps.
Report object presence, color binding, relation correctness, and diversity;
do not select only attractive examples. This is a proposed evaluation
design, not empirical evidence of a model's prompt capability.

\subsection{Provenance and AI assistance}
Figures~\ref{fig:ddpm}, \ref{fig:ldm}, and \ref{fig:attention} are original
TikZ schematics explaining cited methods; no paper figure is reproduced.
Figure~\ref{fig:noise} and Table~\ref{tab:noise} are generated by the supplied
code. Table~\ref{tab:benchmark} is transcribed from LDM Table~6; ratios are
our arithmetic. The two qualitative comparison tables are analytical
syntheses with sources in their captions.

The document uses the requested ICCV 2025 author kit
\cite{template}, with its official style and bibliography files.
It uses 10-point type, two columns, and US Letter paper, following the
new template request. The author block retains HCMUS and CS420 information
and three name/ID placeholders. Template provenance is included in the project.

OpenAI Codex \cite{codex} assisted with drafting, equations, code,
diagram construction, source checking, and LaTeX preparation.
Preparation checks included consulting the original DDPM/LDM equations and
LDM Table~6, recalculating the numerical examples and ratios, and compiling
and visually reviewing the PDF. The named student authors should review
and accept the final manuscript and replace the author placeholders before
submission. No human study, new trained model, or pretrained generation
benchmark is claimed.

\pagecut
\twocolumn[{
\section{Notation, Assumptions, and Additional Derivation}\label{app:notation}
\subsection{Ranges and units}
Images are normalized and dimensionless. Clean pixels in the illustration
lie in $[-1,1]$; noisy states, latent values, and predicted noise are
real-valued and need not stay within that interval.
There is no physical time unit in the diffusion index.

\begin{center}
\begin{minipage}{\textwidth}
\centering\small
\captionsetup{hypcap=false}
\captionof{table}{Notation used in the equations; meanings are also introduced
beside the corresponding equations in the main text.}
\begin{tabularx}{\linewidth}{@{}p{3.3cm}Y@{}}
\toprule
Symbols & Meaning, range, and units\\
\midrule
$x_0,x_t,\hat x_0$ & Clean/noisy/reconstructed image; vector in $\R^D$ or
an equivalent image tensor; normalized intensity units\\
$D,H,W,h,w,c,d_z$ & Positive integer dimensions. $D=3HW$ for RGB,
$d_z=hwc$; $H,W$ count pixels, $h,w$ count latent positions\\
$t,T,j$ & Integer timestep, terminal step, and product index;
$1\le t\le T$, $1\le j\le t$\\
$\beta_t,\alpha_t,\ab_t$ & Dimensionless schedule coefficients;
$0<\beta_t<1$, $\alpha_t=1-\beta_t$, $0<\ab_t\le1$\\
$q,p_\theta,p_{\rm data}$ & Fixed forward transition, learned reverse/generative
distribution, and data distribution, respectively\\
$\Normal(\mu,\Sigma),\Id$ & Gaussian with mean $\mu$ and covariance $\Sigma$;
identity covariance of the relevant dimension\\
$\eps,\eps_z,u,\xi$ & Independent standard-normal draws in their respective
sampling operations; each coordinate has mean $0$ and variance $1$\\
$\sigma_t,\mu_\phi,\sigma_\phi$ & Reverse noise standard deviation
$\sigma_t\ge0$; VAE posterior mean and positive standard deviations\\
$\theta,\phi,\psi,\eta,\omega$ & Trainable parameter arrays of denoiser,
VAE encoder, decoder, condition encoder, and illustrative classifier\\
$\tilde z,z_0,z_t,f,a$ & Raw/clean-scaled/noisy latent;
spatial factor $f\ge1$, with $H,W$ divisible by $f$;
fixed dimensionless scale $a>0$\\
$y,\tau_\eta(y),F,C$ & Condition, encoded condition, U-Net features, and token
embedding matrix. Continuous embeddings are real-valued\\
$n,m,d_u,d_y,d_a,d_v$ & Positive integer counts and widths for cross-attention;
see Section~\ref{sec:attention} for all matrix shapes\\
$Q,K,V,A,O$ & Query, key, value, attention-weight, and output matrices;
$A_{ij}\in[0,1]$ and $\sum_j A_{ij}=1$\\
$s,w,p_{\rm drop},\varnothing$ & CFG scale $s\ge0$, alternative scale $w=s-1$,
dropout probability in $(0,1)$, and null condition\\
$K,r,\mathcal T$ & In the cost model only, $K$ counts updates (not attention
keys), $r$ counts evaluations/update, and $\mathcal T$ is seconds\\
$L_{\rm simple},L_{\rm LDM}$ & Expected summed squared noise-prediction
errors; dimensionless squared normalized-coordinate units\\
$\E,\|\cdot\|_2,\KL$ & Expectation, Euclidean norm, and KL divergence;
natural logarithms make KL and negative log likelihood use nats\\
\bottomrule
\end{tabularx}
\end{minipage}
\end{center}
\vspace{8pt}
}]


\pagecut
\subsection{Why the forward marginal is available in closed form}
Assume independent standard-normal innovations
$\eps^{(1)},\eps^{(2)}$ and let $x_0$ be fixed. Applying
Eq.~\eqref{eq:forward} twice gives
\begin{align*}
x_1&=\sqrt{\alpha_1}x_0+\sqrt{\beta_1}\eps^{(1)},\\
x_2&=\sqrt{\alpha_2\alpha_1}x_0
+\sqrt{\alpha_2\beta_1}\eps^{(1)}
+\sqrt{\beta_2}\eps^{(2)}.
\end{align*}
The two independent noise terms have total variance
$\alpha_2\beta_1+\beta_2
=\alpha_2(1-\alpha_1)+(1-\alpha_2)=1-\alpha_2\alpha_1$.
A sum of independent Gaussians is Gaussian. Repeating this argument yields
Eq.~\eqref{eq:closed}, with mean $\sqrt{\ab_t}x_0$ and covariance
$(1-\ab_t)\Id$ \cite{ho2020ddpm}.

Gaussian innovations are independent across coordinates and steps, but
the noisy images themselves have correlations inherited from $x_0$.
When $\ab_T$ is close to zero, dependence on the starting image is weak
and a standard-normal starting prior is a useful approximation.
A finite schedule need not make it exactly equal to that prior.
The learned Gaussian reverse transition is a modeling approximation
motivated by small forward steps; no finite network is assumed to know
the exact reverse conditional.

\subsection{Why noise prediction is useful without knowing the true noise}
For a fixed $x_t$ and $t$, the vector minimizing expected squared prediction
error is $\E[\eps\mid x_t,t]$ (and, in a conditional model,
$\E[\eps\mid x_t,t,y]$). Thus a trained predictor estimates a statistical
direction shared by plausible clean examples, not an oracle that recovers
every random draw. Its corresponding clean-state estimate is
\begin{equation}
\hat x_0(x_t,t)=
\frac{x_t-\sqrt{1-\ab_t}\eps_\theta(x_t,t)}{\sqrt{\ab_t}}.
\end{equation}
This follows by rearranging Eq.~\eqref{eq:closed}; prediction errors remain
and can be magnified at small $\ab_t$. The reverse sampler uses predictions
over many levels rather than treating one estimate as guaranteed recovery.
The same argument applies in latent space, followed by the learned decoder.

\subsection{Two training stages and their different roles}
Stage one updates encoder $\phi$ and decoder $\psi$ to preserve useful
image information in a regularized representation. Stage two fixes them
and trains denoiser $\theta$ with Eq.~\eqref{eq:ldmloss}. The original LDM
conditional formulation jointly optimizes the condition encoder and
denoiser \cite[Sec.~3.3]{rombach2022ldm}; particular later implementations
may instead keep the text encoder fixed.

The distinction matters: improving the denoiser cannot guarantee exact
recovery of details that the first stage systematically discards.
Conversely, a high-quality autoencoder does not itself provide the
conditional latent distribution needed for generation. The VAE prior
regularizer and the learned diffusion prior play different roles.
